\chapter{設計手法について}
\begin{abstract}

本ドキュメントは超音速衝動タービン翼の設計手法をまとめたものであり，
NACA および NASA の古典的な衝撃波・膨張波理論に基づく設計法を整理し，
数値的に翼形状を生成するための手順を示す。
本ドキュメントに対応するプログラム OpenLaval は，
上面・下面で異なるプラントル・マイヤー角を指定できる
\textbf{非対称翼（asymmetric blade）} に対応している。

非対称翼では，上面と下面で膨張量が異なるため，
それぞれ別の R$^\ast$，別の衝撃波打ち消し曲線，
別の円弧区間が生成される。
本ドキュメントでは，この非対称設計を含めた最新の設計手順を記述する。

\end{abstract}

\section{超音速衝動タービンの設計に当たって}

超音速衝動タービンの設計は，
翼前縁で発生する衝撃波を打ち消すように
上面・下面の壁面形状を構築することから始まる。
必要となる物理量は比較的少なく，
主に以下が入力パラメータとなる：

\begin{itemize}
  \item 入力マッハ数 $M_{\mathrm{in}}$
  \item 出口マッハ数 $M_{\mathrm{out}}$
  \item 比熱比 $\gamma$
  \item 入射角 $\beta_{\mathrm{in}}$
  \item 上面・下面のプラントル・マイヤー角（対称または非対称）
\end{itemize}

OpenLaval では，上面と下面で異なる


\[
(\nu_{l,\mathrm{lower}},\ \nu_{u,\mathrm{lower}}),\qquad
(\nu_{l,\mathrm{upper}},\ \nu_{u,\mathrm{upper}})
\]


を指定できるため，
非対称翼の設計が可能である。
これにより，上面と下面で異なる膨張量を持つ翼形状を構築できる。

\section{設計手順}

設計手順の全体像は以下の通りである：

\begin{enumerate}
  \item 上面・下面それぞれについて衝撃波打ち消し曲線（特性曲線）を生成する。
  \item 各面に対して円弧区間を構築する（R$^\ast$ は面ごとに異なる）。
  \item 上面の直線区間を入射角に合わせて構築し，下面の座標と整合させる。
  \item 上面・下面の差から翼高さを決定し，全体座標を構築する。
  \item 必要に応じて前縁の切り落とし（edge offset, edge delta）を適用する。
\end{enumerate}

非対称翼では 1〜3 の処理が
\textbf{上面と下面で独立に行われる}点が特徴である。

\section{非対称プラントル・マイヤー角}

対称翼では単一の $(\nu_l, \nu_u)$ を用いるが，
非対称翼では以下の 4 つを指定する：



\[
\nu_{l,\mathrm{lower}},\quad
\nu_{u,\mathrm{lower}},\quad
\nu_{l,\mathrm{upper}},\quad
\nu_{u,\mathrm{upper}}
\]



これらは必ず


\[
\nu_{\min} \le \nu_l < \nu_u \le \nu_{\max}
\]


を満たす必要がある。
OpenLaval のバリデーションはこの条件を厳密にチェックする。

\section{衝撃波打ち消し曲線}

衝撃波打ち消し曲線は，特性方程式


\[
\phi(R^\ast,\theta) = \mathrm{const.}
\]


に基づき生成される。
非対称翼では，上面と下面で異なる初期角度を用いるため，
2 本の独立した曲線が生成される。

\section{円弧区間}

特性曲線の出口角度と入口角度を整合させるために，
各面ごとに円弧区間を構築する。
非対称翼では


\[
R^\ast_{\mathrm{upper}} \ne R^\ast_{\mathrm{lower}}
\]


となるため，円弧半径も面ごとに異なる。

\section{直線区間}

上面の直線区間は入射角 $\beta_{\mathrm{in}}$ に合わせて構築し，
下面の座標と整合するように切り落とす。
非対称翼では，上面・下面の曲線形状が異なるため，
直線区間の接続点も面ごとに異なる。

\section{前縁の切り落とし}

OpenLaval では，前縁の切り落としを
\texttt{edge\_delta} と \texttt{edge\_offset} により制御する。
非対称翼でも同じ処理が適用されるが，
上面・下面の曲線が異なるため，
切り落とし点の座標は面ごとに異なる。

\section{プログラムの使い方}

OpenLaval の設定ファイルでは以下を指定できる：

\begin{itemize}
  \item 対称翼：\texttt{vl}, \texttt{vu}
  \item 非対称翼：\texttt{vl\_lower}, \texttt{vl\_upper}, \texttt{vu\_lower}, \texttt{vu\_upper}
  \item 入射角：\texttt{beta\_in}
  \item 前縁パラメータ：\texttt{edge\_delta}, \texttt{edge\_offset}
  \item 補間点数：\texttt{num\_points}
\end{itemize}

CLI では \texttt{--asymmetric} フラグを用いて
非対称パラメータを上書きできる。
